\section{Homotopy type theory}
\label{hott}
We unfortunately do not have enough space here to provide an introduction from
scratch to dependent type theory and homotopy type theory, so we refer the reader
to the reference book for an in depth presentation~\cite{hottbook}. The main
purpose of this section is to fix some terminology and notations for classical
notions.

\paragraph{Universe}
We write $\U$ for the \emph{universe}, \ie the large type of all small types,
which we suppose to be closed under dependent sums and products. We write
$\Pi(x:A).B$ or $(x:A)\to B$ for $\Pi$-types, and $A\to B$ for the case
where~$B$ is non-dependent. Similarly, we write $\Sigma(x:A).B$ for
$\Sigma$-types, and $A\times B$ for the non-dependent version. The two
projections from a $\Sigma$-type are respectively written $\fst$ and $\snd$.

\paragraph{Paths}
Given a type $A:\U$ and two elements $a,b:A$, we write $a=_A b$ for the type of
\emph{identities}, or \emph{paths}, between~$a$ and~$b$: its elements are proofs
of equality between $a$ and $b$. In particular, for any $a:A$, the type $a=a$
contains the term $\refl{a}$ witnessing for reflexivity of
equality.
%
We sometimes write $x\defd t$ to indicate that $x$ and $t$ are equal by
definition.
%
The elimination principle of identities, \aka \emph{path induction} and often
noted~$J$, roughly states that, given $a:A$, in order to show a property
$P:(x:A)\to(a=x)\to\U$ for every $x:A$ and $p:a=x$ it is enough to show it in
the case where $x\defd a$ and $p\defd\refl{a}$.
%
% Given a $a:A$ and a property $P:(x:A)\to(a=x)\to\U$, the canonical map
% \begin{align*}
  % \pa{(x:A)\to(p:a=x)\to P\,x\,p}&\to P\,a\,\refl{a}
  % % f&\mapsto f\,\,\refl{a}
% \end{align*}
% given by evaluation at $\refl{a}$ is an equivalence. In particular, it admits
% an inverse, which encodes the induction principle for paths (\aka \emph{path
% induction}) and is commonly noted~$\J$: in order to show a property~$P$
% depending on a path~$p$ whose target endpoint is free, it is enough to show it
% when the path is $\refl{}$.
%
By path induction, the following can be shown. Given a type~$A$ and a type
family $B:A\to\U$, a path $p:x=y$ in~$A$ induces a function
$\subst Bp:B(x)\to B(y)$ witnessing for the fact that equality is
\emph{substitutive}. As a special case, any path $p:A=B$ between two
types~$A,B:\U$ induces a function $\transport p:A\to B$, called the
\emph{transport} along~$p$, as well as an inverse function
$\transportinv p:B\to A$. Finally, given a function $f:A\to B$, any path $p:x=y$
in~$A$ induces a path $\cong f(p): f(x) = f(y)$ witnessing for the fact that
equality is a \emph{congruence}.

\paragraph{Higher inductive types}
Many functional programming languages allow the definition of inductive types,
which are freely generated by constructors. For instance, the type~$\S0$ of
booleans is generated by two elements (true and false).
%
In the context of homotopy type theory, languages such as cubical Agda feature a
useful generalization of such types, called \emph{higher inductive types}. They allow, in addition to traditional constructors for elements of the type,
constructors for equalities between elements of the type. For instance, the type
corresponding to the circle $\S1$ can be defined as generated by two points $a$
and $b$ and two equalities $p,q:a=b$ between those points.
% \begin{center}
% \begin{minipage}{0.4\linewidth}
% \begin{verbatim}
% data S¹ : Type where
  % a b : S¹
  % p q : a ≡ b
% \end{verbatim}
% \end{minipage}
% \begin{minipage}{0.4\linewidth}
  % \[
    % \begin{tikzpicture}[baseline=(current bounding box.center)]
      % \filldraw (-.5,0) circle (0.05) node[left] {$a$};
      % \filldraw (.5,0) circle (0.05) node[right] {$b$};
      % \draw (0,0) circle (.5);
      % \draw (0,.5) node[below]{$p$};
      % \draw (0,-.5) node[above]{$q$};
    % \end{tikzpicture}
  % \]
% \end{minipage}
% \end{center}
Higher-dimensional spheres~$\S n$ can be defined in a similar way.

\paragraph{Univalence}
\label{equivalence}
A map $f:A\to B$ is an \emph{equivalence} when it admits both a left and a right
inverse. In particular, every isomorphism is an equivalence. We write
$A\equivto B$ for the type of equivalences from~$A$ to~$B$. The identity is
clearly an equivalence and we thus have, by path induction, a canonical map
$(A=B)\to(A\equivto B)$ for every types~$A$ and~$B$: the \emph{univalence axiom}
states that this map is itself an equivalence.
%
In particular, every equivalence $A\equivto B$ induces a path $A=B$.
%
It is known that univalence implies the \emph{function extensionality}
principle~\cite[Section~2.9]{hottbook}: given functions $f,g:A\to B$, if
$f(x)=g(x)$ for any $x:A$ then $f=g$ (and the expected generalization to
dependent function types is also valid).

\paragraph{Homotopy levels}
A type~$A$ is \emph{contractible} when the type $\Sigma(x:A).(y:A)\to(x=y)$ is
inhabited: this means that we have a ``contraction point'' $a_0 : A$, and a
continuous family of paths from $a_0$ to every other point in $A$.
%
A type~$A$ is a \emph{proposition} (\resp a \emph{set}, \resp a
\emph{groupoid}) when $(x=y)$ is contractible (\resp a proposition, \resp a
set) for every $x,y:A$. Intuitively, a contractible type is a point (up to
homotopy), a proposition is a point or is empty, a set is a collection of points
and a groupoid is a space which bears no non-trivial 2-dimensional (or higher)
structure.
% For instance, consider the following three types corresponding to
% spheres of respective dimensions $-1$, $0$, $1$ and $2$:
% \AR{Usually (eg in [HoTT Book, §6.5]) the $-1$-sphere is taken to be the empty
% type, not the contractible type, so that $\Sigma\S{-1}\equivto\S0$.}
% \[
  % \begin{array}{c@{\qquad\qquad\qquad}c@{\qquad\qquad\qquad}c@{\qquad\qquad\qquad}c}
    % % \begin{tikzpicture}
      % % \filldraw (0,0) circle (.05);
      % % \draw (0,-.5);
    % % \end{tikzpicture}
    % &
    % \begin{tikzpicture}
      % \filldraw (0,0) circle (.05);
      % \filldraw (1,0) circle (.05);
      % \draw (.5,-.5);
    % \end{tikzpicture}
    % &
    % \begin{tikzpicture}
      % \filldraw (0,0) circle (.05);
      % \filldraw (1,0) circle (.05);
      % \draw (.5,0) circle (.5);
    % \end{tikzpicture}
    % &
    % % https://tex.stackexchange.com/questions/348609/draw-a-3d-sphere-with-radius-with-tikz
    % \begin{tikzpicture}
      % \shade[ball color = gray!40, opacity = 0.4] (.5,0) circle (.5);
      % \draw (.5,0) circle (.5);
      % \draw (0,0) arc (180:360:.5 and 0.2);
      % \draw[dotted] (1,0) arc (0:180:.5 and 0.2);
      % \filldraw (0,0) circle (.05);
      % \filldraw (1,0) circle (.05);
    % \end{tikzpicture}
    % \\
    % \S{-1}&\S0&\S1&\S2
  % \end{array}
  % \]
% For instance, the sphere $\S{-1}$ (\ie the empty type) is a proposition, the
% sphere $\S0$ (\ie the booleans) is a set and the sphere $\S1$ (\ie the circle)
% is a groupoid, and the sphere $\S2$ is not a groupoid (because there is a
% 2-dimensional ``hole'' in it).
%
We write $\Set$ for the type of sets. Given a type~$A$, we write $\isSet(A)$ (\resp $\isGroupoid(A)$) for the predicate indicating that~$A$ is a set (\resp groupoid).

\paragraph{Truncation}
Given a type $A$, its \emph{propositional truncation} turns it into a
proposition in a universal way. It consists of a type~$\ptrunc{A}$, which is a
proposition, equipped with a map~$\ptrunq-:A\to\ptrunc{A}$ such that, for any
proposition~$B$, the map $(\ptrunc A\to B)\to(A\to B)$ induced by precomposition
by~$\ptrunq-$ is an equivalence. Intuitively, the type $\ptrunc A$ behaves like~$A$,
except that we do not have access to its individual elements: the elimination
principle for propositional truncation states that in order to construct an
element of~$B$ from an element of~$\ptrunc A$, we can only assume that we have
an element of~$A$ if~$B$ itself is a proposition.
%
The \emph{set truncation} $\strunc A$ of a type~$A$ is defined similarly, as the
universal way of turning~$A$ into a set, and we write $\strunq{x}$ for the image
of~$x:A$ in the truncation; and we can similarly define the \emph{groupoid truncation}~$\gtrunc A$.

\paragraph{Fibers}
Given a function $f:A\to B$, we write $\fib fb$ for the type $\Sigma(a:A).(f\,a=b)$, called the \emph{fiber} of~$f$ at $b$. The function $f$ is said to be \emph{surjective} when the type $(b:B)\to\ptrunc{\fib fb}$ is inhabited, \ie when every element of $B$ merely admits a preimage.


% Local Variables:
% TeX-master: "article"
% End:
